\section{Archimedean geometry and the single-row score}\label{sec:arch}

\subsection{Singular points and rates}
Let $L$ be the Picard--Fuchs operator of $f$ in the coordinate $t$ (Section~\ref{sec:sources}). Its singular points in the finite $t$-plane are the images of the cusps and elliptic points of $\Gamma$ together with the zeros and poles of any rational factor in $f$; the latter are apparent. Order the non-apparent singularities by modulus, $0<|t_1|\le|t_2|\le\cdots$. Then $a_n\asymp|t_1|^{-n}$ up to polynomial factors, so $\lambda_1=1/|t_1|$. The companion $b$ solves the inhomogeneous equation and $b-\xi a$, for the Ap\'ery limit $\xi$, is the unique combination continuing analytically past $t_1$; its growth is governed by the next singularity, $\lambda_2=1/|t_2|$. Hence
\[
\text{score}:=\log(1/|\lambda_2|)-k,
\qquad
\xi\notin\mathbb Q\ \text{ if }\ \text{score}>0,
\]
where $k$ is the denominator exponent, $d_n^kb_n\in\mathbb Z$.

When a Fricke involution acts on $t$ by $t\mapsto1/(ct)$ the singular points pair off, $t_1t_2=1/c$, and the recurrence is palindromic: $\lambda_1\lambda_2=c$. For Ap\'ery's two rows $c=\pm1$ and $\lambda_2=1/\lambda_1$; this is the precise sense in which levels $5$ and $6$ are perfect. For Zagier's row E (level $8$, Catalan) the singularities are $t=\tfrac18,\tfrac14$, not reciprocal, and $\lambda_2=4>1$: the ratio $b_n/a_n$ converges ($2^{-n}$) but the linear form diverges. This is the geometric reason every scalar modular Catalan system in the literature fails, and the reason the Catalan work of Section~\ref{sec:tworow} uses two rows.

\subsection{Denominators}\label{sec:denoms}
Each Eichler integration $D^{-1}$ costs at most one factor $d_n$ \cite{ProjDomb}, so generically $k=w+1$. The exception is a \emph{free integration}: for Cooper's $s_7,s_{10},s_{18}$ one has $(n+1)\mid A_n$ \cite{Bogner2013}, hence $d_n^{2}B_n\in\mathbb Z$ instead of $d_n^3$. Sharpness of $k$ was verified for every row in Table~\ref{tab:census} to $n=200$.

\subsection{The census}
% Table 1 -- census of Apery-like lattice rows.
% Generated by lattice/paper_tables/gen_census.gp (+ companion scripts in the same
% directory); numbers transcribed by hand from the gp run recorded in
% lattice/paper_tables/census_raw.txt. See paper/tables/README.md for exact
% regeneration commands and normalisation conventions.
\begin{table}[t]
\centering
\caption{Census of lattice rows: parameters $(a,b,c)$ (or defining recurrence), order $w+1$,
level/group where known, archimedean period, $c=\lambda_1\lambda_2$, characteristic roots
$\lambda_{1,2}$ (4 s.f.; for complex-conjugate pairs, $|\lambda_1|=|\lambda_2|=\sqrt{|c|}$ is
listed once), denominator exponent $k$ (sharp where marked \checkmark, verified $n\le200$),
measured $p$-adic slopes $\sigma_p$ at $n\approx300$ (``--'' = no growing slope; predicted value
$v_p(c)+2\kappa_p$ in parentheses where it differs from the naive $v_p(c)$), denominator growth
rate $\kappa_p$ of $a_n$ where $a_n\notin\mathbb Z$, single-row score $\log(1/|\lambda_2|)-k$,
and budget $\log\lambda_1-k+\sum_p\kappa_p\log p$. Companion normalisation $b_0=0,b_1=1$
throughout; see README for the few rows using a different convention.}
\label{tab:census}
\small
\begin{tabular}{l l c c l r r c c c c c c c c}
\hline
name & $(a,b,c)$/def. & $w{+}1$ & level & period & $c$ & $\lambda_1$ & $\lambda_2$ & $k$ &
$\sigma_2$ & $\sigma_3$ & $\sigma_5$ & $\sigma_7$ & score & budget \\
\hline
Zagier A & $(7,2,-8)$ & 2 & 4/6 & $\zeta(2)/4$ & $-8$ & $8.000$ & $-1.000$ & 2\checkmark & 3(3) & 0 & 0 & 0 & $-2.000$ & $0.079$ \\
Zagier B & $(9,3,27)$ & 2 & 9 & none (complex) & $27$ & \multicolumn{2}{c}{$5.196\pm$} & 2\checkmark & 0 & 3(3) & 0 & 0 & $-3.648$ & $-0.352$ \\
Zagier C & $(10,3,9)$ & 2 & 6/9 & $L(2,\chi_{-3})/2$ & $9$ & $9.000$ & $1.000$ & 2\checkmark & 0 & 2(2) & 0 & 0 & $-2.000$ & $0.197$ \\
Zagier D & $(11,3,-1)$ & 2 & 11 & $\zeta(2)/5$ & $-1$ & $11.09$ & $-0.0902$ & 2\checkmark & 0 & 0 & 0 & 0 & $0.406$ & $0.406$ \\
Zagier E & $(12,4,32)$ & 2 & 8 & $G/2$ (Catalan) & $32$ & $8.000$ & $4.000$ & 2\checkmark & 5(5) & 0 & 0 & 0 & $-3.386$ & $0.079$ \\
Zagier F & $(17,6,72)$ & 2 & 6/8 & $5L(2,\chi_{-3})/8$ & $72$ & $9.000$ & $8.000$ & 2\checkmark & 3(3) & 2(2) & 0 & 0 & $-4.079$ & $0.197$ \\
AZ$(7,3,81)$ & $(7,3,81)$ & 3 & Beauville-IV & none (complex) & $81$ & \multicolumn{2}{c}{$9.000\pm$} & 3\checkmark & 0 & 4(4) & 0 & 0 & n/a & $-0.803$ \\
AZ$(9,3,-27)$ & $(9,3,-27)$ & 3 & -- & $L(3,\chi_{-3})/3$\,$^{\dagger}$ & $-27$ & $19.39$ & $-1.392$ & 3\checkmark & 0 & 3(3) & 0 & 0 & $-3.331$ & $-0.035$ \\
Domb $(10,4,64)$ & $(10,4,64)$ & 3 & 6 & $7\zeta(3)/24$ & $64$ & $16.00$ & $4.000$ & 3\checkmark & 6(6) & 0 & 0 & 0 & $-4.386$ & $-0.227$ \\
$\eta=(11,5,125)$ & $(11,5,125)$ & 3 & order-6 & $\tfrac12L(3,\chi_5)$\,$^{\dagger}$ & $125$ & \multicolumn{2}{c}{$11.18\pm$} & 3\checkmark & 0 & 0 & 3(3) & 0 & n/a & $-0.588$ \\
$T=(12,4,16)$ & $(12,4,16)$ & 3 & 6 & $7\zeta(3)/32$ & $16$ & $23.31$ & $0.6863$ & 3\checkmark & 4(4) & 0 & 0 & 0 & $-2.624$ & $0.149$ \\
Ap\'ery $(17,5,1)$ & $(17,5,1)$ & 3 & 5/6 & $\zeta(3)/6$ & $1$ & $33.97$ & $0.02944$ & 3\checkmark & 0 & 0 & 0 & 0 & $0.526$ & $0.526$ \\
Cooper $s_7$ & $A_1{=}4$, order-3 rec. & 3 & 7 & $\zeta(2)/7$ & $-27$ & $27.00$ & $-1.000$ & 2\checkmark & -- & -- & -- & -- & $-2.000$ & $1.296$ \\
Cooper $s_{10}$ & $A_1{=}2$, order-3 rec. & 3 & 10 & $\zeta(2)/5$ & $-64$ & $16.00$ & $-4.000$ & 2\checkmark & -- & -- & -- & -- & $-3.386$ & $0.773$ \\
Cooper $s_{18}$ & $A_1{=}6$, order-3 rec. & 3 & 18 & $\tfrac12L(2,\chi_{-3})$ & $192$ & $16.00$ & $12.00$ & 2\checkmark & -- & 1(1) & -- & -- & $-4.485$ & $0.773$ \\
Zudilin (Catalan) & $Q_0{=}1,Q_1{=}\tfrac74$, ord.-2 & 2 & -- & $G$ & $1$ & $11.09$ & $0.0902$ & n/a$^{\ddagger}$ & 8(8) & -- & -- & -- & n/a$^{\ddagger}$ & $1.18^{\S}$ \\
Nesterenko $(4,7)$ & Pad\'e det.\ construction & -- & -- & $G$ & --$^{\S}$ & --$^{\S}$ & --$^{\S}$ & --$^{\S}$ & --$^{\S}$ & -- & -- & -- & --$^{\S}$ & --$^{\S}$ \\
$L(f,2)$, wt.\ 3 cusp & Domb curve, $a_1{=}2$ & 2 & 12,$\chi_{-3}$ & $L(f,2)$ & $64$ & $16.00$ & $4.000$ & 2\checkmark & 2$^{\P}$(6) & 0 & 0 & 0 & $-3.386$ & $0.773$ \\
$\zeta(7)$, level 24 & Sym$^6$, order-7 & 7 & 24 & $\tfrac{1463}{13824}\zeta(7)$ & --$^{\parallel}$ & $6.828$ & $4.828$ & ?$^{\parallel}$ & none & none & none & none & ?$^{\parallel}$ & ?$^{\parallel}$ \\
\hline
\end{tabular}

\vspace{2mm}
\footnotesize
$^{\dagger}$Period identification cited from \texttt{consolidation/SLOPE\_CENSUS.md} \S2
(archive source-text identification for AZ$(9,3,-27)$; book \texttt{v8/08c\_chi5\_application.tex}
for $\eta$), not independently re-derived by \texttt{lindep} in this pass.
$^{\ddagger}$The Zudilin row is non-integral ($a_n=Q_m\notin\mathbb Z$); the fixed-$k$ column does not
apply and $\kappa_2=4$ carries the denominator cost instead (verified: $v_2(Q_m)=-4m+2s_2(m)$ to
$m\le40$). Its own-row score is likewise not of the fixed-$k$ form; the budget $+1.18$ is cited
from \texttt{consolidation/THEORY\_NOTES\_03\_lattices.md} \S5 rather than recomputed with a
newly-fixed convention here (flagged in \texttt{README.md}).
$^{\S}$Nesterenko $(4,7)$ and the Zudilin budget: $\lambda_{1,2},k,\sigma_p$ are cited from
\texttt{consolidation/CATALAN\_AUDIT.md} (exact rows verified there: $\log B_n/n\to7.6507$,
$\log|4B_nG-C_n|/n\to-7.3109$, $\log|V_nG-U_n|/n\to E_2{=}14.3931$) and were \emph{not}
independently decomposed into a single $(\lambda_1,\lambda_2,k)$ triple within this task; see
\texttt{README.md} for the honest limitation.
$^{\P}$Measured $\sigma_2=2$ for $L(f,2)$ disagrees with the naive prediction $v_2(64)=6$ quoted
in \texttt{consolidation/SPORADIC\_SEARCH.md}; $a_n$ is confirmed integral ($\kappa_2=0$) to
$n=300$, so the discrepancy is not a denominator effect. The budget $\log\lambda_1-k=0.7726$
(independent of $\sigma_2$) matches \texttt{SPORADIC\_SEARCH.md} exactly. See README.
$^{\parallel}$The order-7 recurrence for the level-24 $\zeta(7)$ row was not fit in closed form
(coefficients not fully available in the source notes); $k$, $c$, score and budget are therefore
left undetermined. $\lambda_{1,2}$ are the paper's own stated values ($4+2\sqrt2$ and its
$1/\sqrt2$-scaled partner), confirmed empirically via the decay rate of $B_n/A_n-\text{target}$
to 33 digits at $n=218$.
\end{table}


Table~\ref{tab:census} records, for every row considered in this paper, the period, $c$, $\lambda_{1,2}$, $k$, the slopes and denominator rates of Section~\ref{sec:padic}, the score, and the \emph{budget} $\log\lambda_1-k+\sum_p\kappa_p\log p$ (the score a row would have if all its $p$-adic resource could be harvested; Section~\ref{sec:tworow}). Two rows have positive score: Ap\'ery's. Three have no archimedean limit at all (complex-conjugate roots: Zagier B, Almkvist--Zudilin $(7,3,81)$ and $(11,5,125)$), yet each has a positive $p$-adic slope and, in B's case, a $p$-adic limit shared with C.

\begin{conjecture}[Scalar barrier]\label{conj:barrier}
For trivial character and $w\ge3$, every integral Hauptmodul system of $\operatorname{Sym}^w$ type on a genus-zero group commensurable with $SL_2(\mathbb Z)$ has negative score.
\end{conjecture}
The evidence is the level-$\le25$ eta-quotient census of \cite{ProjectRace}, in which level $6$ is the unique winner for $w=2$, and the absence of any positive-score $\zeta(5)$ or $\zeta(7)$ system among the many constructed in the project (e.g.\ the level-$12$ $\zeta(5)$ system, whose conditional function has a $\log^5$ singularity \emph{on the unit circle}, so that its linear form decays only like $(\log n)^4/n$). Raising the level crowds cusps into the unit $q$-disc and shrinks $|t_2|$ while the required $k=w+1$ grows.
